% Fichier généré automatiquement par tools/generate_corrections.py.
% Source : CIN/CIN-02-VitesseAcceleration/08_RR3D/08_RR3D.tex
% Statut : complet (2/2 question(s) renseignée(s), 0 reprise(s)).
\begin{corrigebox}[Corrigé extrait de la banque]
\CorrigeQuestion{1}
Le point $C$ peut décrire un tore de grand rayon $R$ et de petit rayon $L$ (surface torique uniquement, pas l'intérieur du tore).
\CorrigeQuestion{2}
On  a  $\vect{AC}=H\vect{j_1}+R\vect{i_1}+L\vect{i_2}$ $=H\vj{0}+R\cos\theta\vi{0}-R\sin\theta\vk{0}+L\cos\varphi\vi{1}+L\sin\varphi\vj{1}$ 
$=H\vj{0}+R\cos\theta\vi{0}-R\sin\theta\vk{0}+L\cos\varphi\left(\cos\theta\vi{0}-\sin\theta\vk{0} \right)+L\sin\varphi\vj{0}$ .

On a donc :
$\left\{
\begin{array}{l}
x_C(t)=R\cos\theta+L\cos\varphi\cos\theta \\
y_C(t)= H+L\sin\varphi\\
z_C(t)=  -R\sin\theta-L\cos\varphi\sin\theta\\
\end{array}
\right.
$ dans le repère $\repere{A}{i_0}{j_0}{k_0}$.
\end{corrigebox}
